%% kapa: exam=2 round=28 year=2017 period=1 subject=appraisal_practice subject_ko=감정평가실무 seq=2 number=2 points=30
%% source: Q-Net 감정평가사 기출문제 (cst00309) artlSeq=5201712 file=raw/1008/5201712/2017년 감정평가사 2차시험 기출문제(감정평가 실무).pdf
%% conversion: kordoc 4.24.51 -> tools/gen2.py (auto); review=proofread-sonnet; flags=none
\begin{longquestion}{【문제 2】}
감정평가사 김○○은 산업단지 내의 염색공장으로 사용되었던 오염 토지에 대하여 시가참고 목적의 감정평가를 의뢰받았다. 관련 법규 및 이론을 참작하고 제시된 자료를 활용하여 다음 물음에 답하시오. (30점) \par
\medskip\noindent 물음 1) 오염 전의 토지가액을 구하시오. (10점) \par
\medskip\noindent 물음 2) 오염 후의 토지가액을 구하시오. (15점) \par
\medskip\noindent 물음 3) 오염된 토지의 스티그마(Stigma) 감정평가 방법을 기술하시오. (5점) \par
<자료 1> 대상 토지의 개요 \par
{\small\centering
\begin{tabular}{|c|c|c|c|c|c|c|}
\hline
\makecell{기호} & \makecell{소재지} & \makecell{지목} & \makecell{면적(m\textsuperscript{2})} & \makecell{용도지역} & \makecell{도로교통} & \makecell{형상\newline{}지세} \\
\hline
\makecell{1} & \makecell{서울특별시\newline{}A구 가동 99} & \makecell{공장용지} & \makecell{9,999} & \makecell{준공업\newline{}지역} & \makecell{중로한면} & \makecell{사다리\newline{}평지} \\
\hline
\end{tabular}\par}

<자료 2> 기본적 사항 \par
1. 감정평가 목적 : 시가참고 \par
2. 기준시점 : 2017.07.01 \par
3. 현장조사 : 2017.03.01～2017.07.01 \par
4. 대상 토지는 2005년부터 산업단지 내에 공업용으로 사용되었고, 토양오염이 발견되어 최근 오염조사 및 정화전문업체가 시료채취를 하여 오염여부를 조사 하였음. 대상 토지는 2010.07.01부터 오염이 시작된 것으로 보이며, 현 상황에서 오염정화에 필요한 기간은 2017.07.01부터 3년이 소요될 것으로 예상됨 \par
대상 토지가 속한 인근지역은 최근 주택 건축이 늘어나고 있으며, 대상 토지 역시 Y주택건설(주)이 주택부지로 분양하기 위하여 2015.07.01에 29,997,000,000원에 매입하였음(종전 건물의 철거비용 150,000,000원은 종전 소유자가 부담). Y주택건설(주)은 대상 토지를 주택부지로 분양하기 위하여 기초공사를 하던 중 2016.07.01에 토양이 오염된 것을 발견하였고 관련 조사가 진행 중임 \par
<자료 3> 오염물질 조사사항 \par
{\small\centering\setlength{\tabcolsep}{3pt}
\begin{tabular}{|>{\centering\arraybackslash}p{0.10\linewidth}|>{\centering\arraybackslash}p{0.43\linewidth}|>{\centering\arraybackslash}p{0.39\linewidth}|}
\hline
구분 & 오염요인 & 조사된 오염물질(단위: mg/kg) \\
\hline
대상 토지\newline{}일부 & 공장운영에 따른 배관 부식과\newline{}오염물질 누출로 추정됨 & 트리클로로에틸렌(TCE) : 75\newline{}테트라클로로에틸렌(PCE) : 50\newline{}톨루엔 : 110\newline{}페놀 : 50\newline{}카드뮴 : 110\newline{}납 : 1,300\newline{}6가크롬 : 80\newline{}비소 : 400\newline{}수은 : 60 \\
\hline
\multicolumn{3}{|p{0.93\linewidth}|}{ㆍ감정평가사 김○○은 오염조사 및 정화 전문업체의 조사 보고서를 검토한\newline{}결과, 대상 토지 일부가「토양환경보전법 시행규칙」제1조의5 관련 별표3\newline{}토양오염우려기준을 상당히 초과하였다고 판단함\newline{}ㆍ향후 오염제거 및 정화공사가 필요하며 이는 합리적이라고 판단함\newline{}ㆍ토양오염의 규모는 2,000m\textsuperscript{3}로 조사됨} \\
\hline
\end{tabular}\par}

<자료 4> 거래사례 자료 \par
{\small\centering\setlength{\tabcolsep}{3pt}
\begin{tabular}{|>{\centering\arraybackslash}p{0.16\linewidth}|*{3}{>{\centering\arraybackslash}p{0.255\linewidth}|}}
\hline
구분 & 사례1 & 사례2 & 사례3 \\
\hline
소재지 & 서울특별시\newline{}A구 가동 97 & 서울특별시\newline{}B구 나동 100 & 서울특별시\newline{}C구 다동 101 \\
\hline
지목 & 공장용지 & 공장용지 & 공장용지 \\
\hline
면적(m\textsuperscript{2}) & 9,000 & 8,000 & 7,500 \\
\hline
이용상황 & 공업용 & 공업용 & 공업용 \\
\hline
도로교통 & 중로한면 & 중로한면 & 중로한면 \\
\hline
형상지세 & 사다리 평지 & 사다리 평지 & 사다리 평지 \\
\hline
거래시점 & 2016.09.23 & 2016.09.14 & 2016.11.06 \\
\hline
거래금액(원) & 15,500,000,000\newline{}(@1,722,000원/m\textsuperscript{2}) & 12,500,000,000\newline{}(@1,562,000원/m\textsuperscript{2}) & 35,000,000,000\newline{}(@4,666,000원/m\textsuperscript{2}) \\
\hline
용도지역 & 준공업지역 & 준공업지역 & 준공업지역 \\
\hline
오염여부 & 오염(TCE, PCE 등,\newline{}1,000m\textsuperscript{3} 정화필요) & 오염(TCE, PCE 등,\newline{}500m\textsuperscript{3} 정화필요) & 토양오염 없음 \\
\hline
\end{tabular}\par}

※ 사례3은 정상적인 거래라고 판단함 \par
<자료 5> 시점수정 자료(지가변동률) \par
{\small\centering
\begin{tabular}{|c|c|c|c|}
\hline
\makecell{구분} & \makecell{A구 공업지역} & \makecell{B구 공업지역} & \makecell{C구 공업지역} \\
\hline
\makecell{2016년 9월} & \makecell{-0.041\%} & \makecell{0.021\%} & \makecell{1.081\%} \\
\hline
\makecell{2016년 10월} & \makecell{-0.042\%} & \makecell{1.085\%} & \makecell{0.752\%} \\
\hline
\makecell{2016년 11월} & \makecell{-0.040\%} & \makecell{0.024\%} & \makecell{0.020\%} \\
\hline
\makecell{2016년 12월} & \makecell{-0.044\%} & \makecell{1.083\%} & \makecell{1.080\%} \\
\hline
\makecell{2017년 1월} & \makecell{1.025\%} & \makecell{-1.022\%} & \makecell{1.500\%} \\
\hline
\makecell{2017년 2월} & \makecell{1.124\%} & \makecell{0.099\%} & \makecell{1.670\%} \\
\hline
\makecell{2017년 3월} & \makecell{2.013\%} & \makecell{0.077\%} & \makecell{1.080\%} \\
\hline
\makecell{2017년 4월} & \makecell{-1.012\%} & \makecell{-0.044\%} & \makecell{1.020\%} \\
\hline
\makecell{2017년 5월} & \makecell{0.051\%} & \makecell{0.022\%} & \makecell{0.750\%} \\
\hline
\end{tabular}\par}

※ 2017년 6월 이후의 지가변동률은 현재 미고시인 상태로 직전 월인 2017년 5월 지가변동률을 연장적용 하기로 함 \par
<자료 6> 기타 참고자료 \par
1. 오염 전의 토지가액은 비교방식을 적용하고, 거래단가를 기준으로 산정함 \par
2. 비교요인표 \par
{\small\centering\setlength{\tabcolsep}{3pt}
\begin{tabular}{|c|c|c|c|c|}
\hline
\makecell{구분} & \makecell{본건} & \makecell{거래사례1} & \makecell{거래사례2} & \makecell{거래사례3} \\
\hline
\makecell{지역요인 비교} & \makecell{100} & \makecell{100} & \makecell{98} & \makecell{115} \\
\hline
\makecell{개별요인 비교} & \makecell{100} & \makecell{95} & \makecell{85} & \makecell{135} \\
\hline
\end{tabular}\par}

※ 요인 비교에서 본건과 사례의 가치형성요인 사항에는 오염에 대한 비교요인은 고려되지 않았음 \par
3. 토양오염 조사비용 자료 \par
: 토양오염의 규모는 2,000m\textsuperscript{3}로 조사되었고, 관련 토양오염 조사비용으로 토양이 오염된 규모를 기준으로 1,000,000원/m\textsuperscript{3}을 2017.07.01에 지급함 \par
4. 정화비용 자료 \par
: 정화방법은 생물학적 처리, 화학적 처리 및 열처리를 복합적으로 적용할 예정이며, 정화기간은 3년이 소요될 것으로 추정되고 연간 정화비용은 600,000원/m\textsuperscript{3}이 소요되며 매년 연말에 지급함 \par
5. 정화공사 기간 중 토지이용제약에 따른 임대료손실 자료 \par
: 임대료 조사사항은 향후 4년 간 시장임대료를 기준으로 보증금 3,000,000,000원, 연간 임대료는 600,000,000원이며, 정화공사 기간 중 임대료손실이 예상되고, 임대와 관련된 지출비용은 미미함 \par
6. 스티그마 자료(오염 전 토지가액을 기준으로 한 가치감소분) \par
{\small\centering
\begin{tabular}{|c|c|c|c|c|}
\hline
\makecell{감가율} & \makecell{오염 전} & \makecell{오염된 상태} & \makecell{정화공사 중} & \makecell{정화공사 후} \\
\hline
\makecell{오염조사 전문업체 보고서 기준} & \makecell{0\%} & \makecell{-30\%} & \makecell{-10\%} & \makecell{-5\%} \\
\hline
\makecell{시장조사 자료} & \makecell{0\%} & \makecell{-20\%} & \makecell{-15\%} & \makecell{-10\%} \\
\hline
\end{tabular}\par}

※ 정화공사 기간은 3년이며, 스티그마의 존속기간은 공사완료 후 1년까지 예상됨 \par
※ 본건 스티그마 금액을 산정하는 경우에는 현재 ‘오염된 상태’의 보고서 및 시장조사 자료의 감가율을 기준으로 각각 산정한 후 평균금액을 적용 \par
7. 이율 자료 \par
1) 보증금은 기간초 지급, 임대료 기간말 지급, 보증금 운용이율은 연 2\% 적용함 \par
2) 시장이자율(할인율) 연 6\%, 화폐의 시간가치 고려함 \par
3) 연복리표(이자율 6\% 기준) \par
{\small\centering
\begin{tabular}{|c|c|c|c|}
\hline
\makecell{기간} & \makecell{일시불 내가계수} & \makecell{연금 내가계수} & \makecell{연금 현가계수} \\
\hline
\makecell{3년} & \makecell{1.191016} & \makecell{3.183600} & \makecell{2.673012} \\
\hline
\makecell{4년} & \makecell{1.262477} & \makecell{4.374616} & \makecell{3.465106} \\
\hline
\end{tabular}\par}

8. 기타 \par
1) 토양오염 이외의 악취 등 가치감소요인은 없는 것으로 봄 \par
2) Y주택건설(주) 대표 장○○은 대상 토지 오염으로 인하여 30,000,000원의 정신적 손실이 발생함 \par
3) 종전 소유자(매도인)의 책임사항은 논외로 함 \par
4) 토지단가는 천원 미만 절사, 물음 1), 물음 2)의 토지가액 및 비용산정 등 금액은 백만원 미만 절사함 \par
\end{longquestion}
